%=========================================================================================
% VITAL -- AIAA AVIATION 2027 extended abstract, Modeling and Simulation Technologies
% AIAA conference template (new-aiaa.cls). Red \tbd{...} marks items still to be filled.
%=========================================================================================
\documentclass[conf]{new-aiaa}
\usepackage[utf8]{inputenc}
\usepackage{textcomp}
\usepackage{graphicx}
\usepackage{amsmath,amssymb}
\usepackage{siunitx}
\usepackage{longtable,tabularx}
\usepackage{xcolor}
\usepackage{pifont}
\setlength\LTleft{0pt}

\newcommand{\tbd}[1]{\textcolor{red}{[#1]}}
\newcommand{\cmark}{\ding{51}}
\newcommand{\xmark}{\ding{55}}
\newcommand{\pmark}{$\sim$}
\newcommand{\vect}[1]{\mathbf{#1}}

\title{VITAL: An Aircraft-Agnostic Flight Simulation Framework for Certification-Oriented Analysis of Novel Aircraft}

\author{Krish Bhatt\footnote{\tbd{Job title, Department, Address, AIAA Member Grade}.}
  \tbd{and co-authors}}
\affil{Virginia Tech, Blacksburg, Virginia, 24061}

\begin{document}
\maketitle

%-----------------------------------------------------------------------------------------
\begin{abstract}
Certification of new aircraft, particularly electric and advanced air mobility designs,
increasingly relies on analysis, yet these aircraft have no flight history from which a
simulation can borrow confidence. A credible simulation must therefore represent the
aircraft as it is built, component by component, and must handle the off-nominal cases on
which certification depends, such as a failed motor or engine. Existing flight dynamics
models describe an aircraft mainly through lumped properties and leave such cases to
custom work. This paper presents VITAL (Virtually Integrated Test and Analysis Lab), a
framework in which an aircraft is assembled from components, each placed on its own
reference axis with its own mass and loads. Mass, center of gravity, and inertia are rolled
up from the components for every configuration, including fuel states, dropped stores, and
failed parts, and every force and moment is carried to the center of gravity by one tested
rigid-body transfer. The aircraft is trimmed and linearized in MATLAB, and its modes are
identified automatically; the trimmed state then initializes a Simulink simulation built on
the Aerospace Blockset six-degree-of-freedom block, which runs in real time and is shown live
in a three-dimensional viewer. VITAL is verified against the NASA Engineering and Safety
Center six-degree-of-freedom check-cases using NASA's own comparison metric
\tbd{headline result}.
\end{abstract}

%-----------------------------------------------------------------------------------------
\section*{Nomenclature}
{\renewcommand\arraystretch{1.0}
\noindent\begin{longtable*}{@{}l @{\quad=\quad} l@{}}
$\vect{A}, \vect{B}$ & state and control matrices of the linear model \\
$\vect{C}_{ba}$ & rotation matrix from axis $a$ to axis $b$ \\
$\vect{F}, \vect{M}$ & force, N, and moment, N$\cdot$m \\
$\vect{f}$ & equations of motion, $\dot{\vect{x}} = \vect{f}(\vect{x}, \vect{u})$ \\
$\vect{g}$ & gravitational acceleration, m/s$^2$ \\
$\vect{I}_3$ & 3-by-3 identity matrix \\
$\vect{J}$ & inertia tensor, kg$\cdot$m$^2$ \\
$m$ & mass, kg \\
$\vect{r}$ & position, m \\
$\vect{u}$ & control vector \\
$\vect{v}$ & velocity relative to the Earth, m/s \\
$\vect{x}$ & state vector \\
$\boldsymbol{\omega}$ & angular velocity, rad/s \\
$\boldsymbol{\Omega}_E$ & Earth rotation rate vector, rad/s \\
\multicolumn{2}{@{}l}{Subscripts}\\
$b$ & body axes \\
$cg$ & center of gravity \\
$e$ & Earth-centered, Earth-fixed axes \\
$i$ & component index \\
\end{longtable*}}

%=========================================================================================
\section{Introduction}
%=========================================================================================
\lettrine{C}{ertification} is how a new aircraft earns permission to fly: the applicant shows
the regulator, requirement by requirement, that the design is safe across its whole
operating envelope. For small airplanes in the United States these requirements are 14 CFR
Part 23, which since Amendment 23-64 states safety objectives in performance-based terms and
lets the applicant propose the means of showing compliance~\cite{part23}. That flexibility
opened the door to electric and hybrid propulsion, distributed propulsion, and the advanced
and urban air mobility vehicles now entering certification. It also moved more of the burden
onto analysis. Many requirements are about the airplane when something has gone wrong: climb
and controllability with one engine inoperative, the minimum control speed with the critical
engine failed, behavior after a control failure. For a conventional twin these cases are well
understood from decades of similar aircraft. For an airplane with fourteen electric motors,
wingtip cruise propellers, and blown high-lift surfaces, there is no such history, and the
first credible look at those cases has to come from a model.

That model has to be more than an aerodynamic database attached to a point mass. A failed
wingtip motor changes the yawing moment through its exact position and thrust line; burning
fuel or draining stored energy, dropping a store, or lowering the gear moves the center of
gravity and changes the inertia; a single pilot command may drive many effectors through a
mixer and actuators with their own limits. Certification-oriented analysis of a novel
aircraft therefore needs an aircraft described at the level of its components, with the
off-nominal configurations defined as precisely as the nominal one, and with analyses
(trim, linearization, mode identification, simulation) that work on any such description.

Widely used flight dynamics tools meet parts of this need (Table~\ref{tab:gap}).
JSBSim~\cite{berndt2004jsbsim} is data-defined and mature: thrusters have individual positions
and orientations, point masses and fuel tanks shift the center of gravity, the flight control
system supports actuators and failures, and it computes linear state-space models. Its mass
model, however, is an empty-aircraft inertia with point masses added, not a build-up from
components; its trim routines act on the conventional pilot controls; and it does not
identify the dynamic modes. YASim, the alternative flight model in FlightGear, infers
aerodynamics from a simplified geometry tuned to match stated performance, and
X-Plane~\cite{xplane} computes forces by blade-element theory over each surface but is closed
and offers no trim or linear analysis; neither uses an oblate, rotating Earth. The MATLAB
Aerospace Blockset and Toolbox~\cite{mathworks_asb} provide the six-degree-of-freedom
equations, linearization, and flying-qualities criteria, but the aircraft's mass, inertia,
and center of gravity are entered as given values rather than derived from its parts.
Finally, credibility itself is rarely demonstrated in public. The NASA Engineering and Safety
Center (NESC) published six-degree-of-freedom check-cases computed by seven simulation tools,
including JSBSim~\cite{nesc2015,jackson2015}, but the results are anonymized and carry no
acceptance criterion, so a new tool has no ready way to show where it stands among them.

\begin{table}[hbt!]
\caption{\label{tab:gap} Capabilities of flight simulation tools as described in their
published documentation (\cmark\ supported, \pmark\ partial, \xmark\ not supported or not
documented)}
\centering
\small
\begin{tabular}{lcccccccccc}
\hline\hline
 & \parbox[c]{1.15cm}{\centering Parts on own axes} & \parbox[c]{1.15cm}{\centering Mass build-up from parts}
 & \parbox[c]{1.3cm}{\centering Off-nominal config-urations} & \parbox[c]{1.15cm}{\centering Mixer and actuators}
 & \parbox[c]{1.1cm}{\centering Trim incl. engine out} & \parbox[c]{1.2cm}{\centering Linear model and modes}
 & \parbox[c]{1.2cm}{\centering Certifi-cation criteria} & \parbox[c]{1.0cm}{\centering Oblate, rotating Earth}
 & \parbox[c]{1.0cm}{\centering NESC check-cases} & \parbox[c]{1.1cm}{\centering Real time, 3-D view} \\\hline
JSBSim             & \pmark & \pmark & \pmark & \cmark & \pmark & \pmark & \xmark & \cmark & \cmark$^a$ & \cmark \\
FlightGear YASim   & \pmark & \pmark & \pmark & \xmark & \xmark & \xmark & \xmark & \xmark & \xmark & \cmark \\
X-Plane            & \cmark & \pmark & \pmark & \pmark & \xmark & \xmark & \xmark & \xmark & \xmark & \cmark \\
Aerospace Blockset & \pmark & \xmark & \xmark & \pmark & \cmark & \pmark & \cmark & \cmark & \xmark & \cmark \\\hline
VITAL              & \cmark & \cmark & \cmark & \cmark & \cmark & \cmark & \cmark & \cmark & \cmark & \cmark \\
\hline\hline
\end{tabular}

\vspace{2pt}
\parbox{0.95\textwidth}{\footnotesize $^a$JSBSim took part in the NESC study; the results are
published anonymously. Partial entries: JSBSim places thrusters and point forces on their own
axes but references aerodynamics to one point, adds point masses and tanks to an
empty-aircraft inertia, supports engine cutoff and actuator failures without configuration
definitions, and produces linear models without mode identification. \tbd{Confirm YASim and
X-Plane entries against their documentation.}}
\end{table}

This paper presents VITAL (Virtually Integrated Test and Analysis Lab), a MATLAB and Simulink
framework built around these needs. The question it addresses is whether a single,
component-level description of an aircraft can support both certification-oriented analysis,
including off-nominal cases, and real-time simulation, with its correctness demonstrated
against public reference data rather than assumed. The contributions are a component-level,
aircraft-agnostic model with mass properties and configurations derived from its parts;
analyses that act on any such model, including trim with failed components and automatic
mode identification; a real-time Simulink simulation and live visualization driven by the
same model; and verification against the NESC check-cases reported in NASA's own metric.

%=========================================================================================
\section{Describing an Aircraft in VITAL}
%=========================================================================================
VITAL contains no code for any particular aircraft. An aircraft is a set of parts, each with
an axis, a mass, and optionally a loads model, together with its control channels and a list
of configurations. Everything that the analyses need is derived from this description, which
is what makes the framework aircraft-agnostic. The examples below use a test vehicle modeled
on the NASA X-57 Maxwell~\cite{x57fs,heeg2019}, built from published values where they exist
and representative values elsewhere: 35 parts, two wingtip cruise propellers, twelve high-lift
propellers, 20 control channels, and five configurations.

\subsection{Parts on Their Own Axes}
Each part sits on its own reference axis, defined by an origin and an orientation relative to
its parent, so that the axes form a tree rooted at the body reference point. In the X-57
vehicle, each half wing, each aileron and flap hinge line, each cruise nacelle, and each of the
fourteen propeller hubs has its own axis; a propeller's thrust line is simply the
$x$-direction of its hub axis. A loads model reports its force and moment on its own axis and
about its own origin, and the framework carries every load to the body axes and the center of
gravity:
\begin{equation}
\label{eq:loads}
\vect{F}_b = \sum_i \vect{C}_{b a_i}\,\vect{F}_i, \qquad
\vect{M}_{cg} = \sum_i \left[\vect{C}_{b a_i}\,\vect{M}_i
  + \left(\vect{r}_i - \vect{r}_{cg}\right) \times \vect{C}_{b a_i}\,\vect{F}_i\right],
\end{equation}
where $\vect{C}_{ba_i}$ is composed along the axis tree and $\vect{r}_i$ is the origin of axis
$a_i$ in body coordinates. A canted motor, an aerodynamic model referenced to a different
point, or a moved center of gravity is a change of data. The same transfer is used by every
analysis and by the simulation.

\subsection{Mass Properties and Configurations}
Each part carries only its own mass, center of gravity, and inertia. For a given
configuration, the framework combines the parts that are present:
\begin{equation}
\label{eq:mass}
m = \sum_i m_i, \qquad
\vect{r}_{cg} = \frac{1}{m}\sum_i m_i \vect{r}_i, \qquad
\vect{J}_{cg} = \sum_i \left[\vect{C}_{b a_i} \vect{J}_i \vect{C}_{b a_i}^{T}
  + m_i \left(\vect{d}_i^T\vect{d}_i\,\vect{I}_3 - \vect{d}_i\vect{d}_i^T\right)\right],
\end{equation}
with $\vect{d}_i = \vect{r}_i - \vect{r}_{cg}$. A configuration states which options are
chosen (gear up or down), how full each tank is, which stores have been released, and which
parts have failed. A failed part keeps its mass but produces no loads, which is the physical
meaning of an inoperative engine, as opposed to one at idle. The X-57 vehicle defines cruise,
takeoff, and landing configurations, a configuration with the left cruise motor failed, and
one with an outboard high-lift motor failed. Each produces its own mass, center of gravity,
and inertia from Eq.~\eqref{eq:mass}, and the same description yields a weight statement and
the center-of-gravity travel across configurations.

\subsection{Controls by Role}
Each control channel is a physical effector, such as a surface or a motor, with travel limits
and a role (pitch, roll, yaw, thrust, lift, flap, trim). A mixer maps pilot or autopilot
commands to channels, $\vect{c} = \vect{M}\,\vect{u}_{\mathrm{cmd}} + \vect{b}$, and an
actuator with a lag, a rate limit, and position stops can be attached to any channel. The
analyses look controls up by role rather than by name. The same trim and control-law code
therefore runs the F-16, whose channels are the elevator, ailerons, rudder, and throttle, and
the X-57 vehicle, whose pitch channel is an all-moving stabilator and whose thrust is shared
among fourteen motors.

\subsection{Aircraft Data}
Aircraft are stored in a JSON format in which every value carries its units and its source,
and formulas are restricted to a safe set so that a data file cannot execute code. Models
published in the ANSI/AIAA S-119 (DAVE-ML) standard~\cite{s119} are imported directly; the
NASA F-16 used for verification enters this way, and its source files are checked against
recorded hashes so that every result traces to the data it came from.

%=========================================================================================
\section{Analysis and Simulation}
%=========================================================================================
\subsection{Trim, Linearization, and Modes}
The compiled aircraft provides its equations of motion as a single function,
$\dot{\vect{x}} = \vect{f}(\vect{x}, \vect{u})$. Trim finds the state and controls that satisfy
\begin{equation}
\label{eq:trim}
\vect{f}(\vect{x}^\ast, \vect{u}^\ast) = \vect{0}
\end{equation}
for a stated flight condition, choosing its unknowns by role, so that an asymmetric case such
as one engine inoperative is trimmed by solving for sideslip or bank together with aileron and
rudder. The result carries an explicit status: trimmed, not trimmable with the reason, not
converged, or outside the range of the aerodynamic data. The linear model is
\begin{equation}
\label{eq:lin}
\vect{A} = \left.\frac{\partial \vect{f}}{\partial \vect{x}}\right|_{\ast}, \qquad
\vect{B} = \left.\frac{\partial \vect{f}}{\partial \vect{u}}\right|_{\ast},
\end{equation}
computed by central differences whose step size is checked for convergence, and the
short-period, phugoid, Dutch-roll, roll, and spiral modes are identified automatically from
the eigenvectors.

\subsection{Six-Degree-of-Freedom Simulation in Simulink}
The trimmed state initializes a Simulink model that VITAL generates from code. The Aerospace
Blockset six-degree-of-freedom Earth-centered quaternion block integrates the motion over an
oblate, rotating Earth:
\begin{equation}
\label{eq:eom}
\begin{aligned}
\dot{\vect{r}}_e &= \vect{v}_e, &
\dot{\vect{v}}_e &= \vect{C}_{eb}\frac{\vect{F}_b}{m} + \vect{g}(\vect{r}_e)
  - 2\,\boldsymbol{\Omega}_E \times \vect{v}_e
  - \boldsymbol{\Omega}_E \times \left(\boldsymbol{\Omega}_E \times \vect{r}_e\right), \\
\vect{J}\,\dot{\boldsymbol{\omega}} &= \vect{M}_{cg} - \boldsymbol{\omega} \times \vect{J}\,\boldsymbol{\omega},
\end{aligned}
\end{equation}
with the attitude carried as a quaternion and $\vect{g}$ the WGS-84 gravitation including the
$J_2$ term. At every integration step the forces and moments of Eq.~\eqref{eq:loads} are taken
from the VITAL aircraft model, so the simulation flies exactly the aircraft that was trimmed
and linearized. The block's conventions were measured rather than assumed; its initial body
rates, for example, are taken relative to the local north-east-down frame, so the inertial
rates specified in the NESC cases are converted before use.

The simulation runs paced to real time and sends one state packet per step to a CesiumJS
viewer, which shows the aircraft over terrain with its flight data and control positions. The
position and attitude drawn in the viewer are computed by the same code that writes the offline
replay files, and every live session is recorded, so a live demonstration can always be traced
back to the run that produced it.

\subsection{Certification Criteria}
The analyses feed a rule engine in which each criterion is a machine-readable record with
its source. MIL-F-8785C~\cite{mil8785c} is implemented, giving a flying-qualities level and
margin for each mode across the envelope. Because a configuration is part of the aircraft
description, the same criteria can be evaluated with a component failed. Part 23 criteria for controllability and climb with one engine inoperative
are the next addition.

%=========================================================================================
\section{Verification}
%=========================================================================================
VITAL is developed with an automated test suite of about \tbd{1,100} tests. Each numerical
check compares a result with published data, an independent tool such as the MATLAB
Aerospace Toolbox, or an analytic solution; for example, the load transfer of
Eq.~\eqref{eq:loads} is checked for power invariance, since the power of a load on a moving
rigid body must be the same about every point. To confirm that the tests detect the errors
that matter, 341 deliberate faults, such as a transposed rotation, a dropped product of
inertia, or a reversed Earth rotation, are inserted into copies of the code, and each must
cause its tests to fail.

The external reference is the NESC set~\cite{nesc2015}: ten cases with spheres and bricks that
exercise the equations of motion, gravitation, atmosphere, and wind, and eight cases with the
NASA F-16, trimmed and under autopilot, including flight around the North Pole and across
the International Date Line. NASA compared its tools with the percent difference
\begin{equation}
\label{eq:pct}
\delta u_{\mathrm{pct}} = 100\,
  \frac{\max_t \left| u(t) - u_{\mathrm{ref}}(t) \right|}
       {\max_{t} \left| u_{\mathrm{ref}}(t) \right|},
\end{equation}
where $u_{\mathrm{ref}}$ is the average of the tools (for wrapped angles, one benchmark tool
with a trigonometric angle difference), applying a small set of documented exclusions. VITAL
uses the same metric and exclusions: NASA's published values are first recomputed from the
NESC data files, and VITAL is then measured against the same reference, without being
included in the average.

%=========================================================================================
\section{Preliminary Results}
%=========================================================================================
\subsection{NESC Check-Cases}
Tables~\ref{tab:case110}--\ref{tab:sig1116} reproduce the four summary tables of the NESC
report with a column added for VITAL. \tbd{Two or three sentences on where VITAL lies relative
to the NASA tools for cases 1--10 and for the F-16 cases, and which signals differ, with the
report's own explanation of the F-16 maneuver transients.} Figure~\ref{fig:case131} shows the
altitude change of case 13.1.

\begin{table}[hbt!]
\caption{\label{tab:case110} Largest differences for check-cases 1--10, by case, percent of
the largest reference value (NASA values from Ref.~\cite{nesc2015}, Table 77)}
\centering
\small
\begin{tabular}{llccccccc}
\hline\hline
Case & Name & SIM 1 & SIM 2 & SIM 3 & SIM 4 & SIM 5 & SIM 6 & VITAL \\\hline
1  & Dragless sphere                     & \tbd{} & \tbd{} & \tbd{} & \tbd{} & \tbd{} & \tbd{} & \tbd{} \\
2  & Dragless tumbling brick             & \tbd{} & \tbd{} & \tbd{} & \tbd{} & \tbd{} & \tbd{} & \tbd{} \\
3  & Tumbling brick with damping         & \tbd{} & \tbd{} & \tbd{} & \tbd{} & \tbd{} & \tbd{} & \tbd{} \\
4  & Ball drop, round non-rotating Earth & \tbd{} & \tbd{} & \tbd{} & \tbd{} & \tbd{} & \tbd{} & \tbd{} \\
5  & Ball drop, round rotating Earth     & \tbd{} & \tbd{} & \tbd{} & \tbd{} & \tbd{} & \tbd{} & \tbd{} \\
6  & Ball drop, WGS-84 Earth             & \tbd{} & \tbd{} & \tbd{} & \tbd{} & \tbd{} & \tbd{} & \tbd{} \\
7  & Ball drop, steady wind              & \tbd{} & \tbd{} & \tbd{} & \tbd{} & \tbd{} & \tbd{} & \tbd{} \\
8  & Ball drop, wind shear               & \tbd{} & \tbd{} & \tbd{} & \tbd{} & \tbd{} & \tbd{} & \tbd{} \\
9  & Eastward ballistic shot             & \tbd{} & \tbd{} & \tbd{} & \tbd{} & \tbd{} & \tbd{} & \tbd{} \\
10 & Northward ballistic shot            & \tbd{} & \tbd{} & \tbd{} & \tbd{} & \tbd{} & \tbd{} & \tbd{} \\\hline
\multicolumn{2}{l}{Largest for each tool} & \tbd{} & \tbd{} & \tbd{} & \tbd{} & \tbd{} & \tbd{} & \tbd{} \\
\hline\hline
\end{tabular}
\end{table}

\begin{table}[hbt!]
\caption{\label{tab:sig110} Largest differences across check-cases 1--10, by signal, percent
of the largest reference value (NASA values from Ref.~\cite{nesc2015}, Table 78)}
\centering
\small
\begin{tabular}{lccccccc}
\hline\hline
Signal & SIM 1 & SIM 2 & SIM 3 & SIM 4 & SIM 5 & SIM 6 & VITAL \\\hline
\tbd{rows as in Table 78} & & & & & & & \\\hline
Largest for each tool & \tbd{} & \tbd{} & \tbd{} & \tbd{} & \tbd{} & \tbd{} & \tbd{} \\
\hline\hline
\end{tabular}
\end{table}

\begin{table}[hbt!]
\caption{\label{tab:case1116} Largest differences for the F-16 check-cases, by case, percent of
the largest reference value (NASA values from Ref.~\cite{nesc2015}, Table 79)}
\centering
\small
\begin{tabular}{llcccc}
\hline\hline
Case & Name & SIM 2 & SIM 4 & SIM 5 & VITAL \\\hline
11   & Subsonic, trimmed            & \tbd{} & \tbd{} & \tbd{} & \tbd{} \\
12   & Supersonic, trimmed          & \tbd{} & \tbd{} & \tbd{} & \tbd{} \\
13.1 & Altitude change              & \tbd{} & \tbd{} & \tbd{} & \tbd{} \\
13.2 & Velocity change              & \tbd{} & \tbd{} & \tbd{} & \tbd{} \\
13.3 & Heading change               & \tbd{} & \tbd{} & \tbd{} & \tbd{} \\
13.4 & Course offset                & \tbd{} & \tbd{} & \tbd{} & \tbd{} \\
15   & Circling the North Pole      & \tbd{} & \tbd{} & \tbd{} & \tbd{} \\
16   & Circling the Equator and IDL & \tbd{} & \tbd{} & \tbd{} & \tbd{} \\\hline
\multicolumn{2}{l}{Largest for each tool} & \tbd{} & \tbd{} & \tbd{} & \tbd{} \\
\hline\hline
\end{tabular}
\end{table}

\begin{table}[hbt!]
\caption{\label{tab:sig1116} Largest differences across the F-16 check-cases, by signal, percent
of the largest reference value (NASA values from Ref.~\cite{nesc2015}, Table 80)}
\centering
\small
\begin{tabular}{lcccc}
\hline\hline
Signal & SIM 2 & SIM 4 & SIM 5 & VITAL \\\hline
\tbd{rows as in Table 80} & & & & \\\hline
Largest for each tool & \tbd{} & \tbd{} & \tbd{} & \tbd{} \\
\hline\hline
\end{tabular}
\end{table}

\begin{figure}[hbt!]
\centering
\begin{minipage}[t]{0.32\textwidth}\centering
  \fbox{\parbox[c][1.5in][c]{0.9\linewidth}{\centering\tbd{Altitude, ft}}}\\[2pt]
  \raggedright a)
\end{minipage}\hfill
\begin{minipage}[t]{0.32\textwidth}\centering
  \fbox{\parbox[c][1.5in][c]{0.9\linewidth}{\centering\tbd{Pitch angle, deg}}}\\[2pt]
  \raggedright b)
\end{minipage}\hfill
\begin{minipage}[t]{0.32\textwidth}\centering
  \fbox{\parbox[c][1.5in][c]{0.9\linewidth}{\centering\tbd{Pitch rate, deg/s}}}\\[2pt]
  \raggedright c)
\end{minipage}
\caption{\label{fig:case131} Check-case 13.1, 100-ft altitude change, NASA tools and VITAL:
a) altitude, b) pitch angle, c) pitch rate.}
\end{figure}

\subsection{F-16 Trim and Modes}
At the case-11 condition (10,013 ft, 565.7 ft/s), the VITAL trim matches the NASA reference
trim to 0.0004 deg in pitch attitude, 0.0002 deg in elevator, and 0.0007 percentage points in
throttle. Table~\ref{tab:modes} lists the modes identified at this condition.

\begin{table}[hbt!]
\caption{\label{tab:modes} Modes of the NASA F-16 model at the case-11 trim condition}
\centering
\begin{tabular}{lccc}
\hline\hline
Mode & Eigenvalue, 1/s & Damping ratio & Frequency, rad/s, or time constant, s \\\hline
Short period & $-1.13 \pm 2.23i$     & 0.45  & 2.50 \\
Phugoid      & $-0.0062 \pm 0.0801i$ & 0.077 & 0.080 \\
Dutch roll   & $-0.389 \pm 3.295i$   & 0.117 & 3.32 \\
Roll         & \tbd{$-2.94$}         & --    & 0.34 \\
Spiral       & \tbd{confirm}         & --    & 99 \\
\hline\hline
\end{tabular}
\end{table}

\subsection{Real-Time Simulation and Visualization}
Flying the F-16 in Simulink and streaming it live, the attitude drawn in the viewer matches the
simulated attitude to $2.4\times10^{-6}$ deg, the delay from the simulation to the screen is
7.9 ms at the 95th percentile, and a 40-s session delivered all 4,001 state packets
(Fig.~\ref{fig:cesium}).

\begin{figure}[hbt!]
\centering
\fbox{\parbox[c][2.0in][c]{0.6\textwidth}{\centering\tbd{Screenshot: live CesiumJS view of the
F-16 over Edwards Air Force Base}}}
\caption{\label{fig:cesium} Live view of the F-16 simulation in CesiumJS.}
\end{figure}

%=========================================================================================
\section{Expected Results of the Final Paper}
%=========================================================================================
The final paper will apply the framework to the X-57 vehicle in its off-nominal
configurations: trim and simulation with a cruise motor and with a high-lift motor failed, the
resulting control margins, and the center-of-gravity and inertia changes across
configurations. It will add Part 23 controllability and climb criteria for one engine
inoperative, pilot-in-the-loop flight in the live viewer, and the aircraft studied in
Ref.~\cite{gould2025falco}, and will complete the NESC comparison with pairwise differences
between VITAL and each NASA tool.

%-----------------------------------------------------------------------------------------
\section*{Acknowledgments}
\tbd{Funding and acknowledgments. AI-use statement as required by AIAA.}

%-----------------------------------------------------------------------------------------
\begin{thebibliography}{99}
\bibitem{part23} Federal Aviation Administration, ``Revision of Airworthiness Standards for
  Normal, Utility, Acrobatic, and Commuter Category Airplanes,'' Final Rule, Amendment 23-64,
  \emph{Federal Register}, Vol. 81, No. 250, 30 Dec. 2016, pp. 96572--96704. \tbd{verify}
\bibitem{berndt2004jsbsim} Berndt, J. S., ``JSBSim: An Open Source Flight Dynamics Model in
  C++,'' AIAA Paper 2004-4923, Aug. 2004. https://doi.org/10.2514/6.2004-4923.
\bibitem{xplane} Laminar Research, ``How X-Plane Works,''
  https://www.x-plane.com/desktop/how-x-plane-works/ [retrieved 10 Nov. 2026].
\bibitem{mathworks_asb} The MathWorks, Inc., \emph{Aerospace Blockset} and \emph{Aerospace
  Toolbox User's Guides}, R2025b, Natick, MA, 2025.
\bibitem{nesc2015} Murri, D. G., Jackson, E. B., Shelton, R. O., \tbd{all authors},
  ``Check-Cases for Verification of 6-Degree-of-Freedom Flight Vehicle Simulations,''
  NASA/TM-2015-218675, Vols. I and II, 2015.
\bibitem{jackson2015} Jackson, E. B., Madden, M. M., Shelton, R. O., \tbd{all authors and
  title}, AIAA Paper 2015-1810, Jan. 2015.
\bibitem{x57fs} NASA Armstrong Flight Research Center, ``X-57 Maxwell,'' Fact Sheet FS-109,
  \tbd{year}.
\bibitem{heeg2019} Heeg, J., \tbd{all authors}, ``Whirl Flutter and the Development of the NASA
  X-57 Maxwell,'' \emph{International Forum on Aeroelasticity and Structural Dynamics}, 2019.
\bibitem{s119} American Institute of Aeronautics and Astronautics, \emph{Flight Dynamics Model
  Exchange Standard}, ANSI/AIAA S-119-2011, Reston, VA, 2011.
\bibitem{mil8785c} \emph{Military Specification: Flying Qualities of Piloted Airplanes},
  MIL-F-8785C, U.S. Department of Defense, 1980.
\bibitem{gould2025falco} Gould, J., Shune, N., Solano Sarmiento, D., and Sarojini, D.,
  ``Performance and Regulatory Assessment Trim Formulations for Conceptual Aircraft Design,''
  AIAA Paper 2025-3624, July 2025. https://doi.org/10.2514/6.2025-3624.
\end{thebibliography}

\end{document}
